\documentclass[10pt,a4paper]{article}

\usepackage[margin=2.3cm]{geometry}
\usepackage[T1]{fontenc}
\usepackage[utf8]{inputenc}
\usepackage{lmodern}
\usepackage{microtype}
\usepackage{amsmath}
\usepackage{booktabs}
\usepackage{array}
\usepackage{graphicx}
\usepackage{hyperref}
\usepackage{xcolor}
\usepackage{enumitem}

\hypersetup{
    colorlinks=true,
    linkcolor=blue!50!black,
    citecolor=blue!50!black,
    urlcolor=blue!50!black
}

\setlength{\parindent}{0pt}
\setlength{\parskip}{0.4em}
\setlist{nosep}
\title{\textbf{Autoencoder for Compression}}
\author{Moritz Zideck\\Computer Vision Course}
\date{\today}

\begin{document}

\maketitle

\begin{abstract}
This report documents the development and evaluation of convolutional
autoencoders for learned compression of satellite imagery. A large,
unlabelled collection of Amazon rainforest patches extracted from Sentinel-2
imagery was used for self-supervised reconstruction training. Three model
families were investigated: a baseline convolutional autoencoder, an
autoencoder with an explicit latent-channel bottleneck, and a residual
autoencoder with uniform latent quantisation and a rate--distortion-inspired
loss. The reconstructed images were subsequently evaluated in a downstream
semantic-segmentation task using a separately labelled Amazon rainforest
dataset. This task-oriented experiment measures whether compression preserves
information that is relevant to a computer-vision model rather than only
visual similarity. The results show that high-quality JPEG and WebP preserve
segmentation performance substantially better than the investigated learned
models. The quantised autoencoder nevertheless provides a well-defined latent
representation with theoretical fixed-length compression ratios between
\(6{:}1\) and \(24{:}1\), depending on the number of latent channels.
\end{abstract}

\section{Motivation and Experimental Objective}

Remote-sensing imagery is spatially large and is commonly acquired repeatedly
over extensive geographical areas. Storage and transmission therefore become
important practical constraints. Conventional codecs such as JPEG and WebP
reduce image size through manually designed transforms, quantisation, and
entropy coding. Learned image compression instead optimises an encoder and
decoder directly from data. The encoder maps an image \(x\) to a latent
representation \(z\), while the decoder reconstructs an approximation
\(\hat{x}\):
\begin{equation}
    z = f_{\theta}(x), \qquad
    \hat{x} = g_{\phi}(z).
\end{equation}

The main objective of this project was to investigate whether such a learned
representation can reduce the amount of image data while retaining information
needed by a semantic-segmentation model. The work was implemented in two
Jupyter notebooks. The first notebook, \texttt{image\_compressor.ipynb},
constructs and trains the compression models. The second notebook,
\texttt{image\_segmentation.ipynb}, fine-tunes a segmentation network and
evaluates it on original, autoencoder-reconstructed, JPEG-compressed, and
WebP-compressed images.

This separation also distinguishes the two learning objectives. Autoencoder
training is self-supervised because the input image itself is the target.
Segmentation training is supervised because every image is paired with a
pixel-level mask. The final evaluation therefore tests transfer across both
datasets and tasks.

\section{Theoretical Background}

\subsection{Autoencoders and Lossy Compression}

An autoencoder consists of an analysis transform \(f_{\theta}\) and a synthesis
transform \(g_{\phi}\). Compression is achieved by restricting the spatial
resolution, channel count, numerical precision, or entropy of the latent
representation. If the bottleneck is sufficiently constrained, the model must
retain image structures that are useful for reconstruction while discarding
less important information.

The simplest reconstruction objective in this work is the mean squared error
\begin{equation}
    \mathcal{L}_{\mathrm{MSE}}
    = \frac{1}{N}\sum_{i=1}^{N}(x_i-\hat{x}_i)^2.
\end{equation}
MSE strongly penalises large pixel deviations but can produce visually smooth
results. The explicit-latent model therefore combines the mean absolute error
and MSE:
\begin{equation}
    \mathcal{L}_{\mathrm{rec}}
    = 0.8\,\lVert x-\hat{x}\rVert_1
    + 0.2\,\mathcal{L}_{\mathrm{MSE}}.
\end{equation}

A practical codec additionally requires discrete symbols. Continuous
floating-point activations cannot be directly represented as a compact bit
stream. The quantised model maps a latent activation from \([-1,1]\) to one of
256 uniformly spaced levels. Rounding is non-differentiable, so a
straight-through estimator is used: the forward pass contains the rounded
value, while the backward pass approximates the derivative by the identity.
The loss is
\begin{equation}
    \mathcal{L}
    = 0.85\,\lVert x-\hat{x}\rVert_1
    + 0.15\,\mathcal{L}_{\mathrm{MSE}}
    + \lambda\,\frac{1}{|z_q|}\sum_i |z_{q,i}|,
\end{equation}
where \(z_q\) is the quantised latent tensor and \(\lambda\) is the rate weight.
The final term is a simple proxy encouraging smaller latent magnitudes. It is
not a learned probability model and does not directly estimate the number of
encoded bits. Modern learned codecs typically optimise an explicit
rate--distortion objective using a probability model and entropy coding
\cite{balle2018}.

\subsection{Semantic Segmentation as a Task-Oriented Quality Measure}

Semantic segmentation assigns a class label to every pixel. Fully
convolutional networks perform this dense prediction without requiring a
separate classifier for each image region \cite{long2015}. In this project, a
pretrained FCN-ResNet50 model was adapted to two output classes and fine-tuned
on Amazon rainforest masks.

The primary evaluation metric is mean intersection over union (mIoU). For
class \(c\),
\begin{equation}
    \mathrm{IoU}_c =
    \frac{\mathrm{TP}_c}
    {\mathrm{TP}_c+\mathrm{FP}_c+\mathrm{FN}_c},
    \qquad
    \mathrm{mIoU} = \frac{1}{C}\sum_{c=1}^{C}\mathrm{IoU}_c.
\end{equation}
The Dice score provides a related overlap measure,
\begin{equation}
    \mathrm{Dice}_c =
    \frac{2\mathrm{TP}_c}
    {2\mathrm{TP}_c+\mathrm{FP}_c+\mathrm{FN}_c}.
\end{equation}
Unlike a reconstruction-only metric, these measures expose whether compression
removes boundaries, textures, or colour differences required by the
segmentation network.

\section{Data Preparation}

\subsection{Self-Supervised Sentinel-2 Patches}

The compression dataset was assembled from Sentinel-2 Level-2A imagery of the
Amazon rainforest. The notebook contains an optional data-preparation pipeline
that searches Microsoft Planetary Computer for scenes with less than
10\,\% cloud cover, stacks bands B04, B03, and B02 into an RGB representation,
and forms a temporal median composite. Random \(512\times512\) patches are then
extracted from geographically larger scenes.

The local experiment data contain 331 training patches and 200 validation
patches. Each stored array has shape \(3\times512\times512\), uses
\texttt{uint16}, and represents Sentinel reflectance values. During loading,
values above 255 are divided by 10,000 and clipped to \([0,1]\). Since no
labels are required, the larger variety of scenes can be used directly for
reconstruction learning.

\subsection{Supervised Amazon Rainforest Segmentation Dataset}

The downstream task uses ``Amazon Rainforest dataset for semantic segmentation
V2'' \cite{bragagnolo2020}. The public dataset is derived from Sentinel-2
Level-2A RGB bands and provides \(512\times512\) images with associated binary
masks. Its published split contains 1,123 training images, 100 validation
images, and 100 test images. The notebook loader pairs images and masks,
converts masks to two classes, and resizes both to \(512\times512\).

For the documented full notebook run, a deterministic subset of 30 training
images and 15 validation images was used. This reduced experiment was
sufficient to verify the complete training and comparison workflow but limits
the statistical strength of the conclusions. The compression and segmentation
datasets are intentionally separate: the autoencoder learns general
reconstruction from the larger unlabelled patch corpus, while the labelled
dataset measures the preservation of task-relevant information.

\section{Implemented Compression Models}

\subsection{Baseline Convolutional Autoencoder}

The baseline architecture contains a configurable sequence of convolution,
batch-normalisation, ReLU, and max-pooling blocks. Every encoder block halves
the spatial dimensions and increases the channel count. The decoder uses
transposed convolutions to restore the image resolution, followed by a sigmoid
output. Eight checkpoints cover encoder depths from one to four and base
channel counts of 32 and 64. This architecture was useful for studying the
effect of depth, but it does not define a low-channel, quantised bottleneck.
Consequently, several internal representations contain more bytes than the
original 8-bit RGB image.

\subsection{Explicit-Latent Autoencoder}

The second architecture uses a convolutional stem, three stride-two
downsampling stages, a \(1\times1\) convolution to an explicit latent width,
and a symmetric decoder. Models were produced with base widths 32 and 64 and
latent widths 8, 16, and 32. A hyperbolic tangent bounds the latent values, and
small uniform noise can be added during training to approximate the
perturbation caused by later quantisation. However, the evaluated checkpoints
store and process the latent tensor as \texttt{float32}; therefore this model
is a compact representation but not yet a complete compressed-file format.

\subsection{Quantised Compressive Autoencoder}

The final architecture adds residual blocks before and after every scale
change. The encoder depth is four, the base width is 32, and two residual
blocks are used at each stage. The latent tensor is quantised to 256 values,
corresponding to an ideal fixed-length representation of eight bits per latent
symbol. Nine models combine latent widths \(L\in\{8,16,32\}\) with rate weights
\(\lambda\in\{10^{-4},10^{-3},5\cdot10^{-3}\}\). They were trained for 40
epochs using Adam with a learning rate of \(10^{-3}\), a training batch size of
4, and a validation batch size of 8.

\section{Compression-Rate Calculation}

The input is treated as an 8-bit RGB image for codec comparison, giving
\begin{equation}
    B_{\mathrm{input}} = 512\cdot512\cdot3\cdot8
    \quad\text{bits}.
\end{equation}
For every model, the latent payload is calculated as
\begin{equation}
    B_{\mathrm{latent}} = H_zW_zC_zb,\qquad
    \mathrm{bpp} = \frac{B_{\mathrm{latent}}}{512\cdot512},\qquad
    R = \frac{B_{\mathrm{input}}}{B_{\mathrm{latent}}},
\end{equation}
where \(H_z,W_z,C_z\) are the latent dimensions and \(b\) is the number of
bits used for each latent value. A ratio above \(1{:}1\) indicates
compression; a ratio below \(1{:}1\) indicates expansion.

\subsection{Baseline-Model Rates}

For a baseline model with encoder depth \(d\) and base width \(B\), the
bottleneck has
\begin{equation}
    C_z=B\,2^{d-1},\qquad H_z=W_z=\frac{512}{2^d}.
\end{equation}
These models retain float32 activations, so \(b=32\). Table~\ref{tab:legacy-rates}
shows that none of the eight baseline checkpoints compresses the raw 8-bit
RGB payload.

\begin{table}[ht]
\centering
\small
\caption{Latent payload calculation for all baseline autoencoders.}
\begin{tabular}{rrrrr}
\toprule
Depth & Base width & Bottleneck shape & Float32 bpp & Input:latent ratio \\
\midrule
1 & 32 & \(32\times256\times256\) & 256 & \(0.094{:}1\) \\
2 & 32 & \(64\times128\times128\) & 128 & \(0.188{:}1\) \\
3 & 32 & \(128\times64\times64\)  & 64  & \(0.375{:}1\) \\
4 & 32 & \(256\times32\times32\)  & 32  & \(0.750{:}1\) \\
1 & 64 & \(64\times256\times256\) & 512 & \(0.047{:}1\) \\
2 & 64 & \(128\times128\times128\)& 256 & \(0.094{:}1\) \\
3 & 64 & \(256\times64\times64\)  & 128 & \(0.188{:}1\) \\
4 & 64 & \(512\times32\times32\)  & 64  & \(0.375{:}1\) \\
\bottomrule
\end{tabular}
\label{tab:legacy-rates}
\end{table}

\subsection{Explicit-Latent and Quantised-Model Rates}

The explicit-latent and quantised architectures perform three spatial
downsampling operations, so the latent shape is \(L\times64\times64\).
Table~\ref{tab:latent-rates} includes every explicit-latent checkpoint. Base
width changes the encoder and decoder capacity but not the latent payload.
Because these checkpoints use float32 latent values, their implemented
representation rates are \(6{:}1\), \(3{:}1\), and \(1.5{:}1\). The 8-bit
column is a hypothetical rate that would require adding quantisation to this
model family.

\begin{table}[ht]
\centering
\small
\caption{Latent payload calculation for all explicit-latent autoencoders.}
\begin{tabular}{rrrrrr}
\toprule
Base & \(L\) & Latent shape & Float32 bpp & Float32 ratio &
Hypothetical 8-bit ratio \\
\midrule
32 & 8  & \(8\times64\times64\)  & 4  & \(6{:}1\)   & \(24{:}1\) \\
32 & 16 & \(16\times64\times64\) & 8  & \(3{:}1\)   & \(12{:}1\) \\
32 & 32 & \(32\times64\times64\) & 16 & \(1.5{:}1\) & \(6{:}1\) \\
64 & 8  & \(8\times64\times64\)  & 4  & \(6{:}1\)   & \(24{:}1\) \\
64 & 16 & \(16\times64\times64\) & 8  & \(3{:}1\)   & \(12{:}1\) \\
64 & 32 & \(32\times64\times64\) & 16 & \(1.5{:}1\) & \(6{:}1\) \\
\bottomrule
\end{tabular}
\label{tab:latent-rates}
\end{table}

All nine quantised checkpoints use 256 levels, or eight bits per latent
symbol. Their complete rate calculation is shown in
Table~\ref{tab:quantised-rates}. The rate weight affects training but does not
change the fixed-length payload for a given \(L\).

\begin{table}[ht]
\centering
\small
\caption{Fixed-length latent rates for all quantised autoencoders.}
\begin{tabular}{rrrrr}
\toprule
\(\lambda\) & \(L\) & Latent shape & bpp & Input:latent ratio \\
\midrule
\(10^{-4}\) & 8  & \(8\times64\times64\)  & 1 & \(24{:}1\) \\
\(10^{-4}\) & 16 & \(16\times64\times64\) & 2 & \(12{:}1\) \\
\(10^{-4}\) & 32 & \(32\times64\times64\) & 4 & \(6{:}1\) \\
\(10^{-3}\) & 8  & \(8\times64\times64\)  & 1 & \(24{:}1\) \\
\(10^{-3}\) & 16 & \(16\times64\times64\) & 2 & \(12{:}1\) \\
\(10^{-3}\) & 32 & \(32\times64\times64\) & 4 & \(6{:}1\) \\
\(0.005\)   & 8  & \(8\times64\times64\)  & 1 & \(24{:}1\) \\
\(0.005\)   & 16 & \(16\times64\times64\) & 2 & \(12{:}1\) \\
\(0.005\)   & 32 & \(32\times64\times64\) & 4 & \(6{:}1\) \\
\bottomrule
\end{tabular}
\label{tab:quantised-rates}
\end{table}

Thus the quantised models theoretically remove 95.8\,\%, 91.7\,\%, and
83.3\,\% of the raw RGB payload for 8, 16, and 32 latent channels,
respectively. All reported rates exclude model parameters because the same
decoder can be reused for many images. They also exclude headers and do not
exploit symbol probabilities. No entropy coder or serialised latent-file
container was implemented, so the tables report representation rates rather
than measured on-disk compression ratios.

\section{Segmentation Experiment}

The segmentation notebook creates an FCN-ResNet50 model with pretrained
torchvision weights and replaces its output for two classes. All parameters
are fine-tuned using AdamW at a learning rate of \(10^{-4}\) for 20 epochs.
The \texttt{VisionTrainer} saves latest and best checkpoints, while
\texttt{LoopEvaluator} and a shared metric accumulator calculate loss, mIoU,
mean accuracy, and Dice score.

For a fair comparison, every image variant uses the same validation samples,
masks, normalisation, segmentation model, and metric implementation. The
baseline receives the original image. Autoencoder variants pass the image
through a loaded checkpoint before normalisation. JPEG and WebP variants are
encoded and decoded in memory at quality settings 35, 65, and 90. Failures are
recorded per variant so that one incompatible checkpoint does not terminate
the complete experiment.

\section{Results and Discussion}

\subsection{Complete Segmentation Results}

Tables~\ref{tab:classic-results}--\ref{tab:quantised-results} report every
completed variant from the full 15-image validation run. The delta is measured
relative to the original-image mIoU of 0.8412.

\begin{table}[ht]
\centering
\small
\caption{Original-image and conventional-codec results.}
\begin{tabular}{lrrr}
\toprule
Input variant & mIoU & Dice & \(\Delta\)mIoU \\
\midrule
Original images & 0.8412 & 0.9137 & 0.0000 \\
JPEG, quality 35 & 0.7125 & 0.8320 & -0.1286 \\
JPEG, quality 65 & 0.7690 & 0.8694 & -0.0722 \\
JPEG, quality 90 & 0.8370 & 0.9113 & -0.0042 \\
WebP, quality 35 & 0.5936 & 0.7439 & -0.2476 \\
WebP, quality 65 & 0.7249 & 0.8404 & -0.1163 \\
WebP, quality 90 & 0.8316 & 0.9080 & -0.0096 \\
\bottomrule
\end{tabular}
\label{tab:classic-results}
\end{table}

\begin{table}[ht]
\centering
\small
\caption{Results for all baseline autoencoders.}
\begin{tabular}{rrrrr}
\toprule
Depth & Base width & mIoU & Dice & \(\Delta\)mIoU \\
\midrule
1 & 32 & 0.6014 & 0.7493 & -0.2398 \\
2 & 32 & 0.5391 & 0.6978 & -0.3021 \\
3 & 32 & 0.3795 & 0.5237 & -0.4616 \\
4 & 32 & 0.3024 & 0.4353 & -0.5388 \\
1 & 64 & 0.6978 & 0.8212 & -0.1434 \\
2 & 64 & 0.4383 & 0.5964 & -0.4029 \\
3 & 64 & 0.4623 & 0.6191 & -0.3789 \\
4 & 64 & 0.5152 & 0.6750 & -0.3260 \\
\bottomrule
\end{tabular}
\label{tab:legacy-results}
\end{table}

\begin{table}[ht]
\centering
\small
\caption{Results for all explicit-latent autoencoders.}
\begin{tabular}{rrrrr}
\toprule
Base width & \(L\) & mIoU & Dice & \(\Delta\)mIoU \\
\midrule
32 & 8  & 0.2767 & 0.3714 & -0.5644 \\
32 & 16 & 0.2680 & 0.3567 & -0.5732 \\
32 & 32 & 0.2611 & 0.3431 & -0.5801 \\
64 & 8  & 0.2593 & 0.3455 & -0.5819 \\
64 & 16 & 0.2612 & 0.3435 & -0.5800 \\
64 & 32 & 0.2620 & 0.3461 & -0.5792 \\
\bottomrule
\end{tabular}
\label{tab:latent-results}
\end{table}

\begin{table}[ht]
\centering
\small
\caption{Results for all quantised autoencoders.}
\begin{tabular}{rrrrr}
\toprule
\(\lambda\) & \(L\) & mIoU & Dice & \(\Delta\)mIoU \\
\midrule
\(10^{-4}\) & 8  & 0.2587 & 0.3420 & -0.5824 \\
\(10^{-4}\) & 16 & 0.4857 & 0.6475 & -0.3554 \\
\(10^{-4}\) & 32 & 0.2724 & 0.3631 & -0.5688 \\
\(10^{-3}\) & 8  & 0.4733 & 0.6317 & -0.3678 \\
\(10^{-3}\) & 16 & 0.2674 & 0.3555 & -0.5737 \\
\(10^{-3}\) & 32 & 0.2726 & 0.3638 & -0.5685 \\
\(0.005\)   & 8  & 0.2632 & 0.3472 & -0.5779 \\
\(0.005\)   & 16 & 0.4028 & 0.5519 & -0.4383 \\
\(0.005\)   & 32 & 0.4641 & 0.6233 & -0.3770 \\
\bottomrule
\end{tabular}
\label{tab:quantised-results}
\end{table}

\clearpage

The original-image baseline reached an mIoU of 0.8412 and a Dice score of
0.9137. JPEG at quality 90 was almost unchanged, losing only 0.0042 mIoU.
WebP at quality 90 also remained close to the baseline. This indicates that
the segmentation model is robust to mild conventional compression.

The strongest learned-model result was obtained by the shallow legacy
autoencoder with 64 base channels. Its mIoU of 0.6978 is substantially below
the baseline and below JPEG quality 35. More importantly, this legacy model
does not provide an actual storage reduction when its float32 bottleneck is
counted. Its accuracy is therefore not evidence of a favourable
rate--distortion trade-off.

Among the genuinely quantised models, \(L=16\) and
\(\lambda=10^{-4}\) performed best with mIoU 0.4857. According to
Table~\ref{tab:quantised-rates}, this model has a theoretical fixed-length ratio of
\(12{:}1\). The \(L=8,\lambda=10^{-3}\) model achieved a similar mIoU of
0.4733 at a more aggressive \(24{:}1\) ratio. The results do not show a
monotonic relationship between latent width, rate weight, and segmentation
accuracy. For example, increasing the latent width did not consistently
improve the downstream task. This suggests that optimisation stability and
the reconstruction objective were more influential than latent capacity
alone.

The explicit-latent, non-quantised models generally performed poorly. Their
reconstructions often caused the segmentation network to favour one class,
which is visible in the very low IoU for the other class in the recorded
per-class metrics. A pixel-reconstruction objective does not explicitly
protect semantic boundaries. Colour shifts or smoothing that appear
acceptable visually can therefore be destructive for a model trained on the
original image distribution.

\section{Comparison with Published Compression Tests}

Published compression results must be compared cautiously because the
datasets, bit-depths, codecs, and quality metrics differ. Standard learned
image-compression papers commonly use natural-image datasets such as Kodak and
report PSNR or MS-SSIM as a function of measured bits per pixel. In contrast,
the present experiment uses Sentinel-2 imagery and reports segmentation mIoU
after reconstruction. The comparison is therefore methodological rather than
a direct leaderboard comparison.

Ball\'e et al.\ implemented a complete nonlinear transform codec containing
quantisation, a learned probability model, and entropy coding
\cite{balle2017}. Their measured rate--distortion curves generally exceeded
JPEG and JPEG 2000, particularly under perceptual MS-SSIM evaluation. Later,
Cheng et al.\ improved the entropy model with discretised Gaussian mixtures
and attention, obtaining performance comparable to the VVC reference codec on
natural-image PSNR tests \cite{cheng2020}. At a constrained rate of
0.15 bpp, a related low-rate model reported MS-SSIM values of 0.9761 and
0.9802 on the CLIC validation and test sets, respectively
\cite{cheng2020lowrate}. These systems differ fundamentally from the current
prototype because they optimise an expected code length and generate an
entropy-coded bitstream. The rate term used in this project only penalises the
mean latent magnitude; it cannot exploit non-uniform symbol frequencies.

L\"ohdefink et al.\ provide a more task-relevant comparison
\cite{lohdefink2019}. They evaluated JPEG, JPEG2000, WebP, and a GAN-based
autoencoder using downstream semantic-segmentation mIoU. At 0.0625 bpp, their
GAN reconstruction exceeded JPEG2000 by 3.9 absolute mIoU percentage points,
even though JPEG2000 had a 5.91 dB PSNR advantage. When the segmentation model
was trained on GAN reconstructions, its mIoU improved by almost 20 absolute
percentage points relative to applying a model trained only on original
images. This result supports two important interpretations of the present
experiment. First, reconstruction PSNR and semantic usefulness need not agree.
Second, the domain mismatch between original and reconstructed images can be
reduced by adapting the downstream model. The current segmentation network was
trained on original images and then tested on autoencoder outputs, so it did
not receive this adaptation benefit.

The closest Earth-observation comparison is the recent study by Molli\`ere
et al.\ \cite{molliere2025}. It compared JPEG2000 with a learned codec using a
discretised Gaussian-mixture likelihood on fire, cloud, and building
segmentation. Learned compression clearly outperformed JPEG2000 on the large
10-channel Sentinel-2 building dataset, but JPEG2000 remained superior on the
smaller single-channel thermal datasets. Their end-to-end compression and
segmentation optimisation also failed to improve over separately optimised
components. For one low-quality building experiment, joint optimisation
changed the result from 0.00008 bpp, 25.76 dB, and F1 0.4925 to 0.4407 bpp,
49.80 dB, and F1 0.7090. A high-quality starting point remained approximately
0.61 bpp with F1 near 0.72. The authors concluded that learned compression is
most advantageous when abundant training data and spectral redundancy are
available.

\begin{table}[ht]
\centering
\small
\caption{Contextual comparison with related compression evaluations. Values
are not directly comparable across datasets and metrics.}
\begin{tabular}{p{3.0cm}p{2.35cm}p{2.25cm}p{6.2cm}}
\toprule
Study & Domain and task & Rate & Main observation \\
\midrule
This project & Sentinel-2 RGB forest segmentation & 1--4 bpp
& Quantised AE mIoU 0.2587--0.4857; JPEG/WebP quality 90 remained close to
the 0.8412 baseline. \\
L\"ohdefink et al. & Automotive semantic segmentation & 0.0625 bpp
& GAN codec exceeded JPEG2000 by 3.9 mIoU points; training segmentation on
decoded images produced a much larger gain. \\
Molli\`ere et al. & EO fire, cloud, and building segmentation & variable;
example 0.4407 bpp
& Learned coding won on the large multispectral dataset, while JPEG2000 was
stronger on smaller single-channel datasets. \\
Cheng et al. & Natural-image reconstruction & 0.15 bpp
& Low-rate learned codec reached 0.9761/0.9802 MS-SSIM on CLIC
validation/test data. \\
\bottomrule
\end{tabular}
\label{tab:literature}
\end{table}

Compared with these studies, the present quantised models use substantially
higher fixed-length rates of 1--4 bpp while producing a larger loss of
downstream accuracy. This does not imply that autoencoder compression is
unsuitable for Sentinel-2 imagery. Rather, it shows that an explicit entropy
model, a measured bitstream, greater training diversity, and decoder-domain
adaptation are central parts of competitive learned compression. The strong
JPEG/WebP results in this test are consistent with Molli\`ere et al.'s finding
that traditional codecs remain difficult to beat when the learned system is
small or insufficiently matched to its data and task.

\section{Limitations and Future Work}

The principal limitation is the size of the downstream experiment: only 30
training and 15 validation images were selected from the larger public
dataset. The values in Tables~\ref{tab:classic-results}--\ref{tab:quantised-results}
demonstrate the implemented
workflow but should not be interpreted as a definitive benchmark. Repeating
the experiment on the full published split and reporting means and confidence
intervals across several random seeds would strengthen the analysis.

The quantised model is also not a complete learned codec. Its magnitude
penalty is only a rate proxy, and the quantised symbols are neither modelled
probabilistically nor entropy-coded. A future implementation should learn a
latent probability model, encode symbols with arithmetic or range coding, and
measure actual byte lengths. This would permit direct rate--distortion curves
against JPEG and WebP in bits per pixel.

Further improvements include training with perceptual or multi-scale losses,
adding an edge-preservation term, and including the segmentation loss during
compressor training. The latter would create a task-aware codec that explicitly
preserves information needed for forest segmentation. Finally, reconstruction
quality should be reported with PSNR and SSIM alongside downstream metrics to
separate visual fidelity from task fidelity.

\section{Conclusion}

The project established an end-to-end experimental pipeline for learned
satellite-image compression and downstream semantic-segmentation evaluation.
Three generations of autoencoder architecture were implemented, culminating
in a residual model with an explicit, uniformly quantised latent bottleneck.
Theoretical fixed-length compression ratios of \(24{:}1\), \(12{:}1\), and
\(6{:}1\) were obtained for latent widths 8, 16, and 32. However, the
segmentation experiment showed a clear quality gap between the learned models
and conventional high-quality JPEG or WebP.

The main conclusion is that dimensional reduction alone is not sufficient for
effective image compression. A useful learned codec requires discrete latent
symbols, an accurate rate model, entropy coding, and an objective aligned with
the information that must be preserved. The notebooks provide the foundation
for this next step and, importantly, evaluate compression according to its
effect on a real computer-vision task.

\begin{thebibliography}{9}

\bibitem{balle2018}
J. Ball\'e, D. Minnen, S. Singh, S. J. Hwang, and N. Johnston,
``Variational Image Compression with a Scale Hyperprior,''
\emph{International Conference on Learning Representations}, 2018.
\url{https://arxiv.org/abs/1802.01436}.

\bibitem{balle2017}
J. Ball\'e, V. Laparra, and E. P. Simoncelli,
``End-to-End Optimized Image Compression,''
\emph{International Conference on Learning Representations}, 2017.
\url{https://arxiv.org/abs/1611.01704}.

\bibitem{cheng2020}
Z. Cheng, H. Sun, M. Takeuchi, and J. Katto,
``Learned Image Compression with Discretized Gaussian Mixture Likelihoods
and Attention Modules,''
\emph{IEEE Conference on Computer Vision and Pattern Recognition}, 2020.
\url{https://arxiv.org/abs/2001.01568}.

\bibitem{cheng2020lowrate}
Z. Cheng, H. Sun, and J. Katto,
``Low Bitrate Image Compression with Discretized Gaussian Mixture
Likelihoods,''
\emph{CVPR Workshop on Learned Image Compression}, 2020.
\url{https://arxiv.org/abs/2004.04318}.

\bibitem{lohdefink2019}
J. L\"ohdefink, A. B\"ar, N. M. Schmidt, F. H\"uger, P. Schlicht, and
T. Fingscheidt,
``GAN- vs. JPEG2000 Image Compression for Distributed Automotive
Perception: Higher Peak SNR Does Not Mean Better Semantic Segmentation,''
2019. \url{https://arxiv.org/abs/1902.04311}.

\bibitem{molliere2025}
C. Molli\`ere, I. Cumplido, M. Zeulner, L. Liesenhoff, M. Schubert, and
J. Gottfriedsen,
``Learned Image Compression for Earth Observation: Implications for
Downstream Segmentation Tasks,'' 2025.
\url{https://arxiv.org/abs/2512.01788}.

\bibitem{long2015}
J. Long, E. Shelhamer, and T. Darrell,
``Fully Convolutional Networks for Semantic Segmentation,''
\emph{IEEE Conference on Computer Vision and Pattern Recognition}, 2015.
\url{https://arxiv.org/abs/1411.4038}.

\bibitem{bragagnolo2020}
L. Bragagnolo, R. V. da Silva, and J. M. V. Grzybowski,
``Amazon Rainforest Dataset for Semantic Segmentation V2,''
Zenodo, 2020.
\url{https://doi.org/10.5281/zenodo.3994970}.

\end{thebibliography}

\end{document}
